\documentclass[11pt,a4paper]{ltjsarticle}
\usepackage[haranoaji]{luatexja-preset}
\usepackage{lmodern}
\usepackage{amsmath,amssymb,bm,mathtools}
\usepackage{geometry}
\usepackage{booktabs,array,tabularx,longtable}
\usepackage{enumitem}
\usepackage{xcolor}
\usepackage{hyperref}
\usepackage{tikz}
\usepackage{fancyvrb}
\usetikzlibrary{arrows.meta,positioning,fit,calc}
\geometry{margin=23mm}
\hypersetup{colorlinks=true,linkcolor=blue!45!black,urlcolor=blue!45!black,
pdftitle={実験I-C・逐次観測方策PoC},pdfauthor={Toru Takahashi}}
\definecolor{Main}{HTML}{176B87}
\definecolor{Accent}{HTML}{EF8354}
\definecolor{Purple}{HTML}{7E57C2}
\definecolor{Soft}{HTML}{EAF4F7}
\definecolor{Warm}{HTML}{FFF3EC}
\setlist[itemize]{leftmargin=1.7em,itemsep=3pt,topsep=4pt}
\setlist[enumerate]{leftmargin=2em,itemsep=4pt,topsep=4pt}
\renewcommand{\baselinestretch}{1.07}
\renewcommand{\arraystretch}{1.18}
\setlength{\parskip}{3pt}
\setlength{\parindent}{1em}
\raggedbottom
\newcommand{\code}[1]{\texttt{#1}}
\newenvironment{takeaway}{\begin{center}\begin{minipage}{0.92\linewidth}\color{Main}\hrule\vspace{6pt}\color{black}}{\vspace{6pt}\color{Main}\hrule\color{black}\end{minipage}\end{center}}
\newenvironment{warningbox}{\begin{center}\begin{minipage}{0.92\linewidth}\color{Accent}\hrule\vspace{6pt}\color{black}}{\vspace{6pt}\color{Accent}\hrule\color{black}\end{minipage}\end{center}}

\title{\textbf{実験I-C・逐次観測方策PoC}\\
\large 一度しかない観測機会に、有限な注意をどう配分するか}
\author{Toru Takahashi}
\date{v0.1 / 2026年8月}

\begin{document}
\maketitle

\begin{takeaway}
\textbf{この実験で確かめること}

実験I-Bでは、観測前の欠落と、観測後の表現圧縮を分けて測った。しかし観測先はOracleが与えていた。本PoCは、毎時刻A/Bの一方だけを見る逐次環境を作り、何を保存したいかに応じて観測方策を学習する。同時刻に見なかった値は後から取得できず、下流の推定器を改良しても完全には戻らない。
\end{takeaway}

推奨環境はGoogle Colabの標準CPUである。GPU、オープンモデル、大容量データ、有料クラウドは使わない。同梱コードはNumPyのみに依存し、full設定も通常のCPUで数十秒以内に終了する。

\tableofcontents

\section{実験I-Bからの進展}

実験I-Bは、観測作用 $\mathcal A_S$ で選ばれなかったチャネルが、後段の表現容量 $C_R$ を増やしても回復しないことを確かめた。ただし、どのチャネルを見るかは、全候補を列挙するOracleが決めていた。

本実験は観測を時間上に展開する。時刻 $t=1,\dots,T$ にA側とB側の値が同時に生じるが、行為主体は一方しか観測できない。

\begin{equation}
a_t\in\{A,B\},\qquad
o_t=x_{t,a_t},\qquad
x_{t,\bar a_t}\notin o_{\leq T}.
\end{equation}

$\bar a_t$ は選ばなかった側を表す。この式の要点は、見なかった同時刻の値が、後の観測履歴の中に存在しないことである。取得後の圧縮ではなく、一度性を持つ観測機会の配分を対象にする。

\begin{center}
\begin{tikzpicture}[node distance=8mm and 8mm,>=Latex,
box/.style={draw=Main,rounded corners,fill=Soft,minimum width=24mm,minimum height=10mm,align=center},
miss/.style={draw=Accent,rounded corners,fill=Warm,minimum width=24mm,minimum height=10mm,align=center}]
\node[box] (world) {時刻$t$\\$A_t\quad B_t$};
\node[box,right=of world] (policy) {観測方策\\$\pi_t(A),\pi_t(B)$};
\node[box,right=of policy] (obs) {選択した値\\$o_t$};
\node[box,right=of obs] (state) {状態推定\\$\hat A,\hat B$};
\node[miss,below=of obs] (lost) {未選択の値\\後から再取得不能};
\draw[->] (world)--(policy);
\draw[->] (policy)--(obs);
\draw[->] (obs)--(state);
\draw[->,Accent] (policy)--(lost);
\end{tikzpicture}
\end{center}

\begin{warningbox}
\textbf{このPoCがまだ示さないこと}

A/Bは、心理志向／客観志向などの意味を持つ軸ではなく、競合する二つの信号源である。方策も現在状態に応じて変わらず、時刻だけに依存する。ここで校正するのは、逐次的な取得機会と観測方策の測定基盤である。
\end{warningbox}

\section{生成データと時間的重要度}

$T=4$ とし、各時刻で相関するA/B信号を生成する。
\begin{align}
A_t&\sim\mathcal N(0,1),\\
B_t&=\rho A_t+\sqrt{1-\rho^2}\,\varepsilon_t,
\qquad \varepsilon_t\sim\mathcal N(0,1),\quad \rho=0.35.
\end{align}

観測値には標準偏差0.05の独立ノイズを加える。A/Bの値自体は各時刻で対称だが、評価で使う時間的重要度は反転させる。

\begin{center}
\begin{tabular}{c|rrrr}
\toprule
& $t=1$ & $t=2$ & $t=3$ & $t=4$ \\
\midrule
Aの重み $w_t^A$ & 1.8 & 1.4 & 0.6 & 0.2 \\
Bの重み $w_t^B$ & 0.2 & 0.6 & 1.4 & 1.8 \\
\bottomrule
\end{tabular}
\end{center}

この設計により、Aを重視するなら前半のA観測が、Bを重視するなら後半のB観測が有利になる。両方を同じ重みで保存する場合は、前半でA、後半でBを見る 「AABB」が候補になる。ただし、コードはこの答えを与えず、方策学習が自力で収束できるかを測る。

推定対象はA/Bの時系列全体である。最適線形読み出しの正規化二乗誤差を時間重みで平均し、
\begin{align}
D_A&=\frac{\sum_t w_t^A\,\mathrm{NMSE}(A_t,\hat A_t)}{\sum_t w_t^A},\\
D_B&=\frac{\sum_t w_t^B\,\mathrm{NMSE}(B_t,\hat B_t)}{\sum_t w_t^B},\\
J_\lambda&=\lambda D_A+(1-\lambda)D_B
\end{align}
とする。$\lambda=1$ はAだけ、$\lambda=0$ はBだけ、$\lambda=0.5$ は両者を同じ重みで保存する条件である。

\section{六つの比較条件}

方策が学習できたかだけでは、何が利得を生んだか分からない。そこで、上限、下限、特化、下流固定、共同更新を分ける。

\begin{longtable}{>{\bfseries}p{0.22\linewidth}p{0.69\linewidth}}
\toprule
条件 & 役割 \\
\midrule
\endhead
full observation & 毎時刻A/Bの両方を見る。時間的な観測制約のない床。\\
Oracle & $2^4=16$通りの固定スケジュールを列挙し、検証データで最良を選ぶ。学習方策の上限対照。\\
random fixed & 各時刻A/Bを等確率で選ぶ。保存要求を用いない基準。\\
A fixed / B fixed & 四時刻すべてで一方を見る。目的が偏った場合の特化対照。\\
learned / frozen & ランダム方策下で作った推定器を固定し、観測方策だけを更新する。方策単独の寄与を測る。\\
learned / joint & 現在の方策分布の下で最適線形推定器を当て直し、方策と下流推定の最小循環を作る。\\
\bottomrule
\end{longtable}

学習方策は時刻ごとの二値softmaxである。
\begin{equation}
\pi_\theta(a_{1:T})=\prod_{t=1}^{T}\pi_{\theta,t}(a_t),
\qquad a_t\in\{A,B\}.
\end{equation}

全16スケジュールを列挙できるため、Monte Carloサンプリングは使わない。目的値の期待値
\begin{equation}
\mathcal J(\theta)=\sum_s\pi_\theta(s)L(s)
\end{equation}
とその方策勾配を厳密に計算し、Adamでロジットを更新する。共同更新条件では、その反復の方策分布に対する最適推定器を十分統計量から解く。

\section{要因計画と工学的ゲート}

\begin{table}[ht]
\centering
\caption{実行条件}
\begin{tabularx}{0.92\linewidth}{>{\bfseries}lX}
\toprule
項目 & 設定 \\
\midrule
時刻数 & 4 \\
相関 $\rho$ & 0.35 \\
観測ノイズ & 0.05 \\
$\lambda$ & 0, 0.25, 0.5, 0.75, 1 \\
quick & 学習2,048、検証1,024、評価1,024、5 seed \\
full & 学習8,192、検証4,096、評価4,096、20 seed \\
必要環境 & CPU、NumPy、メモリ数GB未満 \\
\bottomrule
\end{tabularx}
\end{table}

事前に固定した校正ゲートは次のとおりである。
\begin{enumerate}
\item 全観測条件の平均歪みが0.01未満である。
\item $\lambda=0.5$ でOracle逐次観測の歪みが全観測より0.05以上大きい。
\item 内点 $\lambda\in\{0.25,0.5,0.75\}$ で、共同更新方策のOracleとの差が0.04未満である。
\item 共同更新方策がランダムを0.05以上上回る。
\item 下流固定方策もランダムを0.03以上上回る。
\item 共同更新方策の最頻スケジュールがOracleと80\% 以上一致する。
\item A/B特化対照の優劣が $\lambda$ に応じて反転する。
\end{enumerate}

これらは人間の志向性の統計検定ではない。逐次観測、歪み評価、方策最適化、対照条件が想定した関係で動くかを確かめる工学的ゲートである。

\section{検証済みfull結果}

20 seedのfull設定は全ゲートを通過した。$\lambda=0.5$ の主要結果を表\ref{tab:result}に示す。

\begin{table}[ht]
\centering
\caption{$\lambda=0.5$ の平均歪み（full、低いほど良い）}
\label{tab:result}
\begin{tabular}{lrrr}
\toprule
方法 & $J_{0.5}$ & $D_A$ & $D_B$ \\
\midrule
全観測 & 0.0025 & 0.0025 & 0.0025 \\
Oracle & 0.1775 & 0.1771 & 0.1779 \\
A特化固定 & 0.4402 & 0.0025 & 0.8779 \\
B特化固定 & 0.4405 & 0.8785 & 0.0025 \\
ランダム固定 & 0.4402 & 0.4404 & 0.4401 \\
学習／下流固定 & 0.1775 & 0.1771 & 0.1780 \\
学習／共同更新 & 0.1775 & 0.1771 & 0.1779 \\
\bottomrule
\end{tabular}
\end{table}

全観測とOracle逐次観測の差は0.1750であった。これは、一時刻に一方しか見られないことの機会費用である。観測方策と下流推定をどれだけ最適化しても、同時刻に見なかった専用情報は完全には回復しない。

学習方策は内点でランダム方策を0.2773上回り、Oracleとの差は0.0001であった。最頻スケジュールは全20 seedでOracleと一致した。

\begin{center}
\begin{tabular}{c|ccccc}
\toprule
$\lambda$ & 0 & 0.25 & 0.5 & 0.75 & 1 \\
\midrule
Oracle & BBBB & ABBB & AABB & AAAB & AAAA \\
学習／固定 & BBBB & ABBB & AABB & AAAB & AAAA \\
学習／共同 & BBBB & ABBB & AABB & AAAB & AAAA \\
\bottomrule
\end{tabular}
\end{center}

この結果は、保存要求の重みが変わると、同じ観測環境でも「いつ何を見るか」が系統的に変わることを校正上の予測どおり示している。

\section{共同更新の利得が出なかった意味}

下流固定と共同更新の差は、内点平均で約0.0001に過ぎなかった。これを「共同更新が成功した」とは解釈しない。

本環境では、A/Bの観測チャネルと推定対象が線形に対応している。ランダム方策下で学習した推定器も、選択された各チャネルを既に十分復号できる。そのため、方策が変わっても、下流推定器を当て直す固有利得がほとんどない。

この無差は、後続実験の対照点になる。現在状態によって有用な観測先が変わり、過去の観測分布が表現器の専門化を生む環境では、固定条件と共同更新条件に差が出る可能性がある。次の問いは、どの条件でその差が必要になるかである。

\section{Google Colab実行手順}

\subsection{用意するファイル}

ZIPを展開し、次の二つを手元に置く。
\begin{itemize}
\item \code{colab\_runner.ipynb}
\item \code{sequential\_policy\_poc.py}
\end{itemize}

\subsection{実行}
\begin{enumerate}
\item Google Colabを開き、「アップロード」から \code{colab\_runner.ipynb} を選ぶ。
\item 最初のコードセルを実行し、表示されたファイル選択で \code{sequential\_policy\_poc.py} を選ぶ。
\item 二番目のセルを実行する。既定は \code{--quick} で、数秒〜数十秒で終了する。
\item 実行終了後、自動判定、方策別性能図、観測確率図を確認する。
\item 最後のセルで結果ZIPをダウンロードする。
\end{enumerate}

quickが成功した後、実行セルの \code{--quick} を \code{--full} に変えると20 seedで再確認できる。GPUランタイムに切り替える必要はない。

\subsection{出力の読み方}
\begin{itemize}
\item \code{result\_report.md}：ゲート判定、主要数値、最頻スケジュール。
\item \code{policy\_performance.svg}：$\lambda$ ごとのOracle、ランダム、二つの学習方策の歪み。
\item \code{learned\_policy.svg}：共同更新条件が時刻ごとにAを見る確率。
\item \code{summary\_results.csv}：条件ごとの平均と標準偏差。
\item \code{raw\_results.csv}：seed別の歪み、最頻スケジュール、Oracleとの一致率。
\item \code{policy\_probabilities.csv}：時刻別のA/B観測確率。
\item \code{learning\_history.csv}：反復中の目的値と方策エントロピー。
\end{itemize}

\section{結論と次の実験}

本PoCは、静的な「どのチャネルを残すか」から、時間的な「いつ何を見るか」へ一歩進んだ。保存要求に応じた方策がOracleと同じスケジュールへ収束し、見逃した時刻の情報に明確な機会費用が生じた。

ただし、これは開ループ方策の校正である。次の非自明な段階では、観測前の信念状態や部分観測によって、次に見るべき対象が変わる環境を用いる。そこで、
\begin{enumerate}
\item 現在状態に条件づけられた閉ループ観測方策、
\item Oracle、ランダム、A/B特化の固定方策対照、
\item 下流固定と共同更新の差が生じる条件、
\item 初期の観測履歴が後の方策と表現に残るか
\end{enumerate}
を問う。この段階で初めて、観測方策と世界表現の学習が互いの条件を変える循環に近づく。

\end{document}
